% ============================================================================
% Model Comparison Section — Add to paper via \input{section_model_comparison}
% ============================================================================

\section{Agent Architecture and Information Comparison}
\label{sec:model-comparison}

We benchmark eight decision-making agents spanning four paradigms: mathematical optimization (Oracle, MSSP), heuristic policies (Base-Stock), reinforcement learning (GNN~V3, TSPPO, Residual~RL), and imitation learning (GNN-IL, Specialist DAgger). Table~\ref{tab:info-access} summarizes each agent's information access, and Table~\ref{tab:architecture} details their architectures.

\subsection{Information Structure}

A critical differentiator among agents is their access to demand information. The Oracle solves a linear program with \emph{perfect} future demand knowledge $D_{t+1:T}$, establishing an upper bound. MSSP approximates this by sampling $N$ demand scenarios and solving a stochastic LP. All learning-based agents operate under \emph{partial observability}---they can only observe current inventory state $X_t$, a 5-period demand history $D_{t-5:t}$, and derived features (demand velocity, time encoding). Crucially, no RL/IL agent has access to future demand.

\begin{table}[h!]
\caption{Information access comparison across agents. \cmark\ = available, \xmark\ = unavailable.}
\label{tab:info-access}
\centering
\small
\begin{tabular}{lccccccc}
\toprule
\textbf{Feature} & \textbf{Oracle} & \textbf{MSSP} & \textbf{Heur.} & \textbf{GNN V3} & \textbf{GNN-IL} & \textbf{DAgger} & \textbf{TSPPO} \\
\midrule
Current inventory $X_t$      & \cmark & \cmark & \cmark & \cmark & \cmark & \cmark & \cmark \\
Pipeline inventory           & \cmark & \cmark & \cmark & via $\text{IP}_t$ & via $\text{IP}_t$ & via $\text{IP}_t$ & via $\text{IP}_t$ \\
\textbf{Future demand} $D_{t+1:T}$ & \cmark & sampled & \xmark & \xmark & \xmark & \xmark & \xmark \\
Demand history $D_{t-5:t}$   & ---    & ---    & \xmark & \cmark & \cmark & \cmark & \cmark \\
Network topology (edges)     & \cmark & \cmark & partial & \cmark & \cmark & \cmark & positional \\
Goodwill sentiment $\psi_t$  & \xmark & \xmark & \xmark & \cmark & \cmark & \cmark & \cmark \\
\bottomrule
\end{tabular}
\end{table}

\subsection{Agent Architectures}

\begin{table}[h!]
\caption{Architectural comparison of benchmark agents.}
\label{tab:architecture}
\centering
\small
\begin{tabular}{lllll}
\toprule
\textbf{Agent} & \textbf{Architecture} & \textbf{Training} & \textbf{Parameters} & \textbf{Decision Process} \\
\midrule
Oracle & LP solver (CBC) & None & 0 & Full-horizon LP with exact $D$ \\
MSSP & Stochastic LP & None & 0 & LP over sampled demand scenarios \\
Heuristic & Analytical & None & 0 & $q = \max(0, \mu(L{+}1) - \text{IP})$ \\
\midrule
GNN V3 & MPNN + attention & PPO (RL) & 600K & Message passing $\to$ pooling $\to$ MLP \\
TSPPO & Transformer encoder & PPO (RL) & 600K & Self-attention over node tokens $\to$ MLP \\
Residual RL & MLP on heuristic & PPO (RL) & 200K & $a = a_{\text{heur}} + f_\theta(\text{obs})$ \\
\midrule
GNN-IL & MPNN + attention & BC $\to$ PPO & 600K & V3 initialized by Oracle imitation \\
Spec.\ DAgger & MPNN + attention & DAgger & 600K$\times$4 & Iterative BC with on-policy Oracle queries \\
\bottomrule
\end{tabular}
\end{table}

\subsection{Training Paradigms}

We identify three training paradigms with distinct characteristics:

\paragraph{Pure Reinforcement Learning (GNN~V3, TSPPO).}
The agent learns entirely from environment interaction via PPO~\cite{schulman2017ppo}. The GNN~V3 uses a Message Passing Neural Network (MPNN) with attention-weighted aggregation and GRU updates over the supply chain graph. TSPPO replaces the GNN with a Transformer encoder~\cite{vaswani2017attention}, treating each supply chain node as a token with learned positional and node-type embeddings. Despite the Transformer's capacity for all-to-all attention, both architectures converge to identical performance ($\approx$82\% of Oracle), suggesting the bottleneck is the \emph{information gap} (no future demand), not the \emph{architecture}.

\paragraph{Imitation Learning (GNN-IL, DAgger).}
The agent learns by imitating Oracle demonstrations rather than exploring from scratch. GNN-IL uses behavioral cloning (BC) to initialize the policy via supervised MSE loss on $(s_t, a^*_t)$ pairs from Oracle rollouts, followed by PPO fine-tuning. This yields 89\% of Oracle performance but suffers from \emph{distribution shift}: the BC policy encounters states during evaluation that differ from Oracle's training trajectory.

Specialist DAgger~\cite{ross2011dagger} addresses this by iteratively: (1) rolling out the \emph{current} policy, (2) querying the Oracle ``what would you do at this state?'', and (3) retraining on the aggregated dataset. This teaches the agent to recover from \emph{its own} mistakes. Training separate DAgger models per demand type yields 92--97\% of Oracle, matching MSSP while additionally handling endogenous goodwill dynamics.

\paragraph{Residual Learning (Residual~RL).}
The agent learns small corrections on top of a base-stock heuristic: $a_t = a_{\text{heur}}(s_t) + f_\theta(s_t)$. This inherits the heuristic's strong baseline performance ($\sim$93\%) and learns refinements, achieving $\sim$92\% of Oracle on tested scenarios.

\subsection{Performance Summary}

\begin{table}[h!]
\caption{Agent performance as percentage of Oracle profit on deterministic demand scenarios (no goodwill), averaged across 5 unseen evaluation seeds. Best RL/IL result per scenario in \textbf{bold}.}
\label{tab:performance}
\centering
\small
\begin{tabular}{lcccccc}
\toprule
\textbf{Scenario} & \textbf{MSSP} & \textbf{Heur.} & \textbf{GNN V3} & \textbf{GNN-IL} & \textbf{Spec.\ DAgger} & \textbf{TSPPO} \\
\midrule
Stationary & 93\% & 98\% & 79\% & 85\% & \textbf{97\%} & 82\% \\
Trend & 97\% & --- & 93\% & 91\% & \textbf{93\%} & --- \\
Seasonal & 93\% & --- & 76\% & 84\% & \textbf{92\%} & --- \\
Shock & 97\% & 87\% & 92\% & 95\% & \textbf{97\%} & --- \\
\midrule
\textbf{Mean} & \textbf{95\%} & $\sim$93\% & 85\% & 89\% & \textbf{95\%} & $\sim$82\% \\
\bottomrule
\end{tabular}
\end{table}

Specialist DAgger matches MSSP's mean performance (95\% of Oracle) while offering two advantages: (i)~sub-millisecond inference time versus MSSP's per-step LP solve, and (ii)~the ability to handle endogenous goodwill dynamics where all learning-based agents \emph{surpass} the Oracle bound by 10--30\%, since the Oracle's LP formulation cannot model sentiment feedback loops.

\subsection{Key Findings}

\begin{enumerate}
    \item \textbf{Information dominates architecture.} GNN~V3 and TSPPO achieve identical performance despite fundamentally different architectures (message passing vs.\ self-attention), confirming that the $\sim$15\% gap to Oracle stems from \emph{missing future demand information}, not model expressiveness.
    
    \item \textbf{Expert demonstrations close the gap.} Imitation learning from Oracle demonstrations (DAgger) recovers most of the performance lost to information asymmetry, raising RL's 82\% to 95\% of Oracle without requiring future demand access at inference time.
    
    \item \textbf{Specialization matters for IL, not RL.} Training demand-type-specific DAgger models improves performance from 76\% (generalist) to 95\% (specialist), while V3 generalist RL already achieves 85\% without specialization.
    
    \item \textbf{RL agents uniquely handle endogenous dynamics.} On goodwill scenarios, all RL/IL agents surpass the Oracle bound, as the Oracle's deterministic LP cannot model the feedback loop between service quality and future demand.
\end{enumerate}
